\begin{frame}{이 과목이 다루는 질문}
  2012년 ImageNet 대회에서 \kb{AlexNet} 이 등장하기 전까지, ``이 사진에
  고양이가 있는가''를 알려주는 방법은 \textbf{사람이 직접 특징을 설계}하는
  것이었다 --- 가장자리 검출기, 색 히스토그램, 특정 모양의 필터를 손으로
  조합해 ``고양이스러움''을 수치로 정의하려 했다.
  \vskip0.6em
  AlexNet은 그런 손설계 없이, 수백만 장의 사진과 정답 라벨만 보고
  \textbf{스스로 무엇을 봐야 하는지 학습}해서 기존 방법들을 압도적인 차이로
  이겼다.
  \vskip1em
  \begin{center}
  \large\kq{규칙을 사람이 짜는 대신, 정답이 있는 데이터로부터\\
      기계가 직접 규칙을 찾아내게 하려면 어떻게 해야 할까?}
  \end{center}
\end{frame}
